\documentclass[11pt,a4paper]{article}
\usepackage[T1]{fontenc}
\usepackage[utf8]{inputenc}
\usepackage{lmodern,amsmath,amssymb}
\usepackage[margin=25mm]{geometry}
\usepackage[hidelinks]{hyperref}
\hypersetup{pdftitle={The large-field region in the Hamiltonian route needs an operator inequality, and the meas},pdfauthor={navstokgap research working notes}}
\newcommand{\tightlist}{\setlength{\itemsep}{0pt}\setlength{\parskip}{0pt}}
\setlength{\emergencystretch}{3em}
\setcounter{secnumdepth}{0}
\begin{document}
\hypertarget{the-large-field-region-in-the-hamiltonian-route-needs-an-operator-inequality-and-the-measure-statement-is-weaker-than-that}{%
\section{The large-field region in the Hamiltonian route needs an
operator inequality, and the measure statement is weaker than
that}\label{the-large-field-region-in-the-hamiltonian-route-needs-an-operator-inequality-and-the-measure-statement-is-weaker-than-that}}

Working out what the gap inequality requires from the large-field region
replaces ``an effective description'' by something more precise, and
corrects one earlier reading. Three positive facts first. \textbf{The
truncation inside the large-field region costs only a longer range.} The
error of truncating the flow-conjugated Hamiltonian at range \(Ka\) is
\(C\exp[2t\|G\|_\infty-cK]\), so on configurations with
\(t\|G\|_\infty\le\eta_\infty/8\) a range
\(K\ge\eta_\infty/(4c)+\log(C/\delta)/c\) brings it below \(\delta\)
there too; the competition recorded in
\href{ground-state-measure-transfer.md}{the transfer note} §4 between
\(e^{-c\eta^2/g^2}\) and \(e^{2\eta}\) was a fixed-\(K\) artifact, and
in the window \(\eta\le C_2\) of Section 4 the range is \(O(1)\).
\textbf{Localizing to the small-field region is cheap per doubling.} For
a smooth partition \(\chi_P^2+\chi_Q^2=1\) depending on the flowed block
density at radius \(2a\), the IMS localization formula gives
\(H=\chi_PH\chi_P+\chi_QH\chi_Q-L\) with
\[L\ \le\ C_L\,\frac{\hbar c}{a}\,\frac{g^2}{\eta^2},\qquad C_L=\frac{32\pi^2c_0^2}{3},\]
supported on the transition region, against the \(\hbar c/a\) scale of
the step. \textbf{A positive perturbation that is large only on a rare
set is harmless}: for \(W\ge0\),
\(\operatorname{gap}(H+W)\ge\operatorname{gap}(H)-\langle\Omega,W\Omega\rangle\),
and the expectation is the measure of the support times the size. The
obstruction is the sign. What the chain of lower bounds needs on the
large-field region is the operator inequality
\[H-E_0\ \ge\ c\,\frac{\hbar c}{a}\,\frac{\eta^2}{g^2}\ \mathbf 1_Q\ -\ (\text{small})\,\mathbf 1_P ,\]
a local energy excess, and every route to it from the Hamiltonian alone
loses a term of order \((\hbar c/a)\times(\text{surface of the block})\)
on the small-field region, because cutting the vacuum costs zero-point
energy on every shared link. That loss is a \emph{negative} perturbation
of order \(\hbar c/a\) on \(P\), and negative perturbations of that size
bind: a well of depth \(D\) on a set of ground-state measure
\(\varepsilon\) captures a state once \(D\) exceeds the localization
cost \(L\), whatever \(\varepsilon\) is. So the Euclidean estimate,
which is a measure statement, is strictly weaker than what the
Hamiltonian chain consumes, which is an operator statement. The polymer
expansion is the device that makes measure statements sufficient, by
attaching the factor \(e^{-c\eta^2/g^2}\) to each large-field polymer as
a convergence factor rather than asking for an operator bound. That is
the reason the constructive programme is Euclidean, stated as a
computation. Constants explicit; nothing promoted.

\hypertarget{what-the-gap-inequality-consumes}{%
\subsection{1. What the gap inequality
consumes}\label{what-the-gap-inequality-consumes}}

The route proves \(\operatorname{gap}(H)\ge\Delta\) by a chain of
operator lower bounds: conjugation (exact), truncation
(\(H\ge(1-\delta)H_{\rm eff}\) in the relative sense, since the
truncated kinetic metric differs from the full one by a small multiple
of it), decimation of the modes between \(a\) and \(2a\) (Feshbach,
\href{weak-coupling-feshbach-reduction.md}{Lemma 1\('\)}), and
induction. The comparison with \(\Delta\) is made only at the end, at
strong coupling, because at any intermediate scale
\(\Delta=\hbar c\Lambda\) is smaller than every error the step can
afford, all of which are measured against \(\hbar c/a\). So each step
must deliver
\[H^{(a)}\ \ge\ H^{(2a)}_{\rm eff}\ -\ R,\qquad R\ \text{small against }\hbar c/a\ \text{and irrelevant},\]
and the question for the large-field region \(Q\) is what it contributes
to \(R\).

\hypertarget{localization-the-ims-formula-on-the-configuration-manifold}{%
\subsection{2. Localization: the IMS formula on the configuration
manifold}\label{localization-the-ims-formula-on-the-configuration-manifold}}

For \(H=K(-\Delta)+V\) on a Riemannian manifold and smooth
\(\chi_1^2+\chi_2^2=1\), the IMS formula (Cycon, Froese, Kirsch and
Simon, \emph{Schrödinger Operators}, Springer 1987, Theorem 3.2;
metadata level) is the identity
\[H=\chi_1H\chi_1+\chi_2H\chi_2-K\big(|\nabla\chi_1|^2+|\nabla\chi_2|^2\big).\]
On \(G^{\mathcal E}\) with \(K=(\hbar c/a)(g^2/2)\) per link, take
\(\chi_P=\cos\theta(f_B)\), \(\chi_Q=\sin\theta(f_B)\) with \(\theta\)
rising from \(0\) to \(\pi/2\) as the flowed block density
\(f_B=\ell\int_B|B_t|^2d^3x\) rises from \(\eta^2\) to \(4\eta^2\), at
\(\ell=2a\) and \(\sqrt{8t}=\ell\). Then
\(|\nabla\chi_P|^2+|\nabla\chi_Q|^2=\theta'(f_B)^2|\nabla f_B|^2\) with
\(\theta'=\pi/(6\eta^2)\), and
\[\frac{\partial f_B}{\partial\theta_{\rm link}}=2\ell a\int_BB_t(x)\,K_t(x-y)\,d^3x\ \le\ 4c_0\,\frac{\eta a}{\ell},\]
since a unit change of a link angle changes the flowed field by
\(aK_t(x-y)\) and \(|B_t|\le2c_0\eta/\ell^2\) on the transition region.
Summing over the \(3(2\ell/a)^3\) links of the smearing neighbourhood,
\(|\nabla f_B|^2\le384c_0^2\eta^2\,\ell/a=768c_0^2\eta^2\) at
\(\ell=2a\), so
\[L_B\ \le\ \frac{\hbar c}{a}\frac{g^2}2\cdot\frac{\pi^2}{36\eta^4}\cdot768c_0^2\eta^2
=\frac{32\pi^2c_0^2}{3}\,\frac{\hbar c}{a}\,\frac{g^2}{\eta^2}.\] Two
remarks. The general-\(\ell\) form carries a factor \(\ell/a\), so
localizing at a scale far above the lattice is expensive and the
localization must be done one doubling at a time, in the variables of
the current scale. And \(L_B\) is a multiplication operator supported on
the transition region of block \(B\), whose ground-state measure is
\(e^{-c\eta^2/g^2}\): it is small against \(\hbar c/a\) by
\(C_Lg^2/\eta^2\) and rare.

\hypertarget{the-sign-of-a-rare-perturbation}{%
\subsection{3. The sign of a rare
perturbation}\label{the-sign-of-a-rare-perturbation}}

\textbf{Proposition 1.} If \(W\ge0\) then
\(\operatorname{gap}(H+W)\ge\operatorname{gap}(H)-\langle\Omega,W\Omega\rangle\).

\emph{Proof.} \(E_1(H+W)\ge E_1(H)\) by min-max, and
\(E_0(H+W)\le\langle\Omega,(H+W)\Omega\rangle=E_0+\langle\Omega,W\Omega\rangle\).
\(\square\)

So a positive perturbation of size \(D\) on a set of ground-state
measure \(\varepsilon\) costs the gap at most \(D\varepsilon\), and with
\(\varepsilon=e^{-c\eta^2/g^2}\) and \(D=O(\hbar c/a)\) that is
negligible.

\textbf{Proposition 2.} A negative perturbation \(-W\) with
\(W\ge D\,\mathbf 1_A\) on a set \(A\) of ground-state measure
\(\varepsilon\) produces a state of energy at most \(E_0+L_A-D\), where
\(L_A=K\int|\nabla\chi_A|^2\Omega^2/\|\chi_A\Omega\|^2\) is the
localization cost of a cutoff \(\chi_A\) supported in \(A\).

\emph{Proof.} The trial state \(\chi_A\Omega/\|\chi_A\Omega\|\) has
energy \(E_0+L_A\) in \(H\) by the exact identity
\(\langle\chi\Omega,(H-E_0)\chi\Omega\rangle=K\int|\nabla\chi|^2\Omega^2\),
and \(-W\) lowers it by at least \(D\). \(\square\)

The measure \(\varepsilon\) does not appear in the conclusion. A
negative perturbation of order \(\hbar c/a\) on a rare set reorganizes
the low spectrum as soon as \(D>L_A\), and \(L_A\) is of order
\((\hbar c/a)C_Lg^2/\eta^2\ll\hbar c/a\). \textbf{The sign of the error
on the large-field region decides everything, and the measure of the
region decides nothing.}

\hypertarget{the-operator-inequality-the-chain-needs-and-its-cost}{%
\subsection{4. The operator inequality the chain needs, and its
cost}\label{the-operator-inequality-the-chain-needs-and-its-cost}}

For the large-field sector to contribute only harmlessly to \(R\), the
step needs
\[\chi_Q(H-E_0)\chi_Q\ \ge\ c_1\,\frac{\hbar c}{a}\,\frac{\eta^2}{g^2}\,\chi_Q^2 ,\]
a \emph{local energy excess}: configurations with a large flowed field
in a block have energy above the vacuum by the field's energy. This is
physically evident at weak coupling and is the Hamiltonian form of the
action lower bound of \href{large-field-action-lower-bound.md}{the
lower-bound note}. Its proof from the Hamiltonian splits \(H=H_N+H'\)
with \(N\) the neighbourhood of the block and uses two facts:
\(V_N\ge(\hbar c/a)(2/g^2)c_1\eta^2\) on \(Q\), by three-dimensional
flow monotonicity and positivity, exactly as in the Euclidean case; and
an upper bound on \(E_0-E_0(H-V_N)\), the energy the neighbourhood's
magnetic terms add to the vacuum. The second is where the loss occurs.
The crude bound uses the ground state of \(H-V_N\), which is Haar-flat
on the interior links and gives
\(\langle V_N\rangle_{\rm Haar}=(\hbar c/a)(2N/g^2)|N|_p\), of the same
order \(1/g^2\) as the excess and useless. The correct order is the
zero-point energy, \(O(1)\cdot(\hbar c/a)|N|_p\), and any trial state
that achieves it must extend the outside vacuum smoothly into \(N\); the
cut then costs \(O(\hbar c/a)\) per shared link, an amount
\[C_V\,\frac{\hbar c}{a}\,|\partial N|\] that is subtracted
\emph{everywhere}, including on \(P\). The resulting inequality is
\[H-E_0\ \ge\ \Big[\frac{2c_1\eta^2}{g^2}-C_V|N|_p\Big]\frac{\hbar c}{a}\,\mathbf 1_Q
\ -\ C_V|\partial N|\,\frac{\hbar c}{a}\,\mathbf 1_P ,\] and the second
term is a negative perturbation of order \(\hbar c/a\) on the
small-field region, of exactly the kind Proposition 2 says cannot be
tolerated. The window in which the first term is positive,
\[\eta\ \ge\ C_1g,\qquad C_1^2=\frac{C_V|N|_p}{2c_1}\ \simeq\ 160\ \text{for a }4a\text{ neighbourhood},\]
together with the small-angle requirement of the perturbative decimation
on \(P\), \(\eta/4\le C_2\), is nonempty only for \(g\le4C_2/C_1\), that
is \(1/g^2\gtrsim10^2\) with these crude constants. Even inside the
window the \(P\)-side loss remains.

\hypertarget{why-the-euclidean-framework-escapes-this}{%
\subsection{5. Why the Euclidean framework escapes
this}\label{why-the-euclidean-framework-escapes-this}}

The Euclidean side proves a measure statement,
\(\mathbb P(Q)\le e^{-c\eta^2/g^2}\), and never an operator one. It can
afford this because the cluster (polymer) expansion consumes measure
statements directly: a large-field polymer enters the expansion with the
weight \(e^{-c\eta^2/g^2}\) as a convergence factor, and the effective
action outside the polymers is computed perturbatively. No inequality of
the form \(H-E_0\ge(\cdots)\mathbf 1_Q\) is ever needed, because there
is no Hamiltonian to bound; the gap is read off at the end from the
decay of the two-point function, which is again a measure statement.

The Hamiltonian chain of Section 1, by contrast, consumes operator
inequalities, and Section 3 shows it cannot substitute a measure
statement for them when the sign is wrong. This is the precise content
of ``the constructive programme is Euclidean'': \textbf{the large-field
region is controllable in measure and uncontrollable in norm, and the
polymer expansion is the device that makes measure control sufficient.}

\hypertarget{consequence-for-state}{%
\subsection{6. Consequence for STATE}\label{consequence-for-state}}

The obligation ``an effective description inside the large-field
region'' is replaced by the following. The truncation there is cheap
(range \(O(\eta)\)), the localization is cheap (\(C_Lg^2/\eta^2\) per
doubling), and positive rare errors are harmless. The chain needs the
local energy-excess operator inequality, whose Hamiltonian proof loses
\(C_V|\partial N|\,\hbar c/a\) on the small-field region, a negative
perturbation of the order that binds. So the Hamiltonian route, as a
chain of operator lower bounds, cannot close through the large-field
region, and the Euclidean route closes it by consuming measure
statements in a polymer expansion. The honest division of labour is:
Hamiltonian for T1, T2, the small-volume theorem, the upper bounds and
the final gap extraction; Euclidean polymer expansion for the
renormalization steps. The coupling map is then: weak-coupling expansion
for \(1/g^2\gtrsim C_1^2/(16C_2^2)\), strong-coupling theorem for
\(54/g^4\le\beta_*\), and the intermediate region between them, in which
the gap forms and no expansion applies; its width in one-loop doublings
is about \(10^3\) with the crude constants above and about \(20\) if the
weak side reaches \(g^2\sim1/2\)
(\href{strong-coupling-threshold-explicit.md}{threshold note} §4).
\end{document}
